% Part: first-order-logic
% Chapter: completeness
% Section: compactness-direct

\documentclass[../../../include/open-logic-section]{subfiles}

\begin{document}

\iftag{FOL}
      {\olfileid{fol}{com}{cpd}}
      {\olfileid{pl}{com}{cpd}}

\olsection{برهان مباشر لمبرهنة التراص}

للتراص برهان مباشر لا يستند إلى مبرهنة الاكتمال، وإن أخذ بطريقة
برهانها. فقد بدأنا هناك بمجموعة متسقة~$\OLId{greek-variable}{Gamma}$ من ال!!{sentence}s،
ومددناها إلى مجموعة متسقة\iftag{FOL}{، مشبعة،}{}
و!!{complete}~$\OLId{greek-variable}{Gamma}^*$ من ال!!{sentence}s. ثم أثبتنا صدق
!!{sentence}s~$\OLId{greek-variable}{Gamma}$ كلها في
\iftag{FOL}{نموذج الحدود~$\Struct{M(\Gamma^*)}$}{!!{valuation}~$\pAssign{\OLId{latin-ordinary}{v}(\OLId{greek-variable}{Gamma}^*)}$}
المنشأ من~$\OLId{greek-variable}{Gamma}^*$، فثبت أن $\OLId{greek-variable}{Gamma}$ تقبل الإرضاء.

والطريق نفسه يثبت أن كل مجموعة من ال!!{sentence}s تقبل الإرضاء
منتهيًا تقبله مطلقًا. وإنما يلزم أن نعيد النتائج السابقة للمّة الصدق
بفرض «قابلة للإرضاء منتهيًا» في موضع فرض «متسقة».

\begin{prop}
\ollabel{prop:fsat-ccs}
افترض أن $\OLId{greek-variable}{Gamma}$ مجموعة من ال!!{sentence}s، !!{complete} وقابلة للإرضاء
منتهيًا. عندئذ:
\begin{enumerate}
\tagitem{prvAnd}{$(!A \land !B) \in \OLId{greek-variable}{Gamma}$
  إذا وفقط إذا كان كل من $!A \in \OLId{greek-variable}{Gamma}$ و $!B \in \OLId{greek-variable}{Gamma}$.}{}

\tagitem{prvOr}{$(!A \lor !B) \in \OLId{greek-variable}{Gamma}$ إذا وفقط إذا كان إما
  $!A \in \OLId{greek-variable}{Gamma}$ وإما $!B \in \OLId{greek-variable}{Gamma}$.}{}

\tagitem{prvIf}{$(!A \lif !B) \in \OLId{greek-variable}{Gamma}$ إذا وفقط إذا كان إما
  $!A \notin \OLId{greek-variable}{Gamma}$ وإما $!B \in \OLId{greek-variable}{Gamma}$.}{}
\end{enumerate}
\end{prop}

\tagprob{FOL}
\begin{prob}
أثبت \olref[fol][com][cpd]{prop:fsat-ccs}. وتجنب استعمال $\Proves$.
\end{prob}
\tagendprob

\tagprob{notFOL}
\begin{prob}
أثبت \olref[pl][com][cpd]{prop:fsat-ccs}. وتجنب استعمال $\Proves$.
\end{prob}
\tagendprob

\iftag{FOL}{%
\begin{lem}
\ollabel{lem:fsat-henkin} يمكن توسيع كل مجموعة من ال!!{sentence}s
قابلة للإرضاء منتهيًا~$\OLId{greek-variable}{Gamma}$ إلى مجموعة مشبعة قابلة للإرضاء
منتهيًا~$\OLId{greek-variable}{Gamma}'$.
\end{lem}
}{}

\tagprob{FOL}
\begin{prob}
برهن \olref[fol][com][cpd]{lem:fsat-henkin}. (إرشاد: مدار البرهان
على أن قبول $\OLId{greek-variable}{Gamma}_\OLId{latin-ordinary}{n}$ الإرضاء منتهيًا يستتبع قبول
$\OLId{greek-variable}{Gamma}_\OLId{latin-ordinary}{n} \cup \{!D_\OLId{latin-ordinary}{n}\}$ له؛ وأثبت ذلك من غير توسل بال!!{derivation}s
ولا بالاتساق.)
\end{prob}
\tagendprob

\iftag{FOL}{
\begin{prop}\ollabel{prop:fsat-instances}
افترض أن $\OLId{greek-variable}{Gamma}$ مجموعة من ال!!{sentence}s كاملة، وقابلة للإرضاء
منتهيًا، ومشبعة.
\begin{tagenumerate}{prvEx,prvAll}
\tagitem{prvEx}{$\lexists[\OLId{latin-ordinary}{x}][!A(\OLId{latin-ordinary}{x})] \in \OLId{greek-variable}{Gamma}$ إذا وفقط إذا كان
  $!A(\OLId{latin-ordinary}{t}) \in \OLId{greek-variable}{Gamma}$ لبعض حد مغلق~$\OLId{latin-ordinary}{t}$ على الأقل.}{}
\tagitem{prvAll}{$\lforall[\OLId{latin-ordinary}{x}][!A(\OLId{latin-ordinary}{x})] \in \OLId{greek-variable}{Gamma}$ إذا وفقط إذا كان
  $!A(\OLId{latin-ordinary}{t}) \in \OLId{greek-variable}{Gamma}$ لكل حد مغلق~$\OLId{latin-ordinary}{t}$.}{}
\end{tagenumerate}
\end{prop}
}{}

\tagprob{FOL}
\begin{prob}
أثبت \olref[fol][com][cpd]{prop:fsat-instances}.
\end{prob}
\tagendprob

\begin{lem}
\ollabel{lem:fsat-lindenbaum} يمكن توسيع كل مجموعة من ال!!{sentence}s
قابلة للإرضاء منتهيًا~$\OLId{greek-variable}{Gamma}$ إلى مجموعة~$\OLId{greek-variable}{Gamma}^*$ !!a{complete}
وقابلة للإرضاء منتهيًا.
\end{lem}

\tagprob{FOL}
\begin{prob}
برهن \olref[fol][com][cpd]{lem:fsat-lindenbaum}. (إرشاد: أثبت أولًا
أن $\OLId{greek-variable}{Gamma}_\OLId{latin-ordinary}{n}$ إذا قبلت الإرضاء منتهيًا، قبلته على الأقل إحدى
المجموعتين $\OLId{greek-variable}{Gamma}_\OLId{latin-ordinary}{n} \cup \{!A_\OLId{latin-ordinary}{n}\}$ و
$\OLId{greek-variable}{Gamma}_\OLId{latin-ordinary}{n} \cup \{\lnot !A_\OLId{latin-ordinary}{n}\}$.)
\end{prob}
\tagendprob

\tagprob{notFOL}
\begin{prob}
برهن \olref[pl][com][cpd]{lem:fsat-lindenbaum}. (إرشاد: أثبت أولًا
أن $\OLId{greek-variable}{Gamma}_\OLId{latin-ordinary}{n}$ إذا قبلت الإرضاء منتهيًا، قبلته على الأقل إحدى
المجموعتين $\OLId{greek-variable}{Gamma}_\OLId{latin-ordinary}{n} \cup \{!A_\OLId{latin-ordinary}{n}\}$ و
$\OLId{greek-variable}{Gamma}_\OLId{latin-ordinary}{n} \cup \{\lnot !A_\OLId{latin-ordinary}{n}\}$.)
\end{prob}
\tagendprob

\begin{thm}[التراص]
\ollabel{thm:compactness-direct} لتكن $\OLId{greek-variable}{Gamma}$ مجموعة من ال!!{sentence}s.
تكون $\OLId{greek-variable}{Gamma}$ قابلة للإرضاء إذا وفقط إذا كانت قابلة للإرضاء منتهيًا.
\end{thm}

\begin{proof}
أما إذا كانت $\OLId{greek-variable}{Gamma}$ قابلة للإرضاء، فهناك
\iftag{FOL}{!!a{structure}~$\Struct{M}$}{!!a{valuation}~$\pAssign{\OLId{latin-ordinary}{v}}$}
تحقق $\iftag{FOL}{\Sat{\OLId{latin-looped}{M}}}{\pSat{\OLId{latin-ordinary}{v}}}{!A}$ لكل $!A \in \OLId{greek-variable}{Gamma}$.
وهذه $\iftag{FOL}{\Struct{M}}{\pAssign{\OLId{latin-ordinary}{v}}}$ ترضي كل جزء منتهٍ
من~$\OLId{greek-variable}{Gamma}$ أيضًا، فالمجموعة $\OLId{greek-variable}{Gamma}$ قابلة للإرضاء منتهيًا.

وأما العكس، فلتكن $\OLId{greek-variable}{Gamma}$ قابلة للإرضاء منتهيًا.\iftag{FOL}{
  نأخذ من \olref{lem:fsat-henkin} مجموعة مشبعة قابلة للإرضاء
  منتهيًا $\OLId{greek-variable}{Gamma}' \supseteq \OLId{greek-variable}{Gamma}$.}{} وتتيح
\olref{lem:fsat-lindenbaum} مدَّ $\iftag{FOL}{\OLId{greek-variable}{Gamma}'}{\OLId{greek-variable}{Gamma}}$
إلى مجموعة~$\OLId{greek-variable}{Gamma}^*$ !!a{complete} تقبل الإرضاء
منتهيًا\iftag{FOL}{، ويبقى إشباع $\OLId{greek-variable}{Gamma}^*$ محفوظًا}{}.
وننشئ \iftag{FOL}{نموذج
  الحدود~$\Struct{M(\Gamma^*)}$}{!!{valuation}~$\pAssign{\OLId{latin-ordinary}{v}(\OLId{greek-variable}{Gamma}^*)}$}
على نحو \olref[mod]{defn:termmodel}.\iftag{FOL}{ فإن
\olref[mod]{prop:quant-termmodel} لم تتوقف على اتساق $\OLId{greek-variable}{Gamma}^*$،
ولا على كونها !!{complete} أو مشبعة؛ وإنما اعتمدت على كون
$\Struct{M(\Gamma^*)}$ مغطاة.}{} ويصح برهان لمّة الصدق
(\olref[mod]{lem:truth}) بعد إحلال
\olref{prop:fsat-ccs} محل \olref[ccs]{prop:ccs}
\iftag{FOL}{، وإحلال
\olref{prop:fsat-instances} محل \olref[hen]{prop:saturated-instances}}{}.
\iftag{FOL}{وأما إذا اشتملت $\OLId{greek-variable}{Gamma}$ على~$\eq$، فنعرّف~$\approx$
بين الحدود المغلقة ونأخذ خارج نموذج الحدود كما في
\olref[ide]{defn:term-model-factor}. وخواص
\olref[ide]{prop:approx-equiv} باقية إذا جعلنا قابلية الإرضاء
منتهيًا مكان الاتساق: فإن الإخفاق في أي من الانعكاسية أو التماثل أو
التعدية أو التوافق مع الدوال والمحمولات تكشفه مجموعة جزئية منتهية
لا تقبل الإرضاء. وعليه يجري برهان حالة الهوية من لمّة الصدق كما في
\olref[ide]{lem:truth}. فالنموذج الناتج يرضي كل ما في~$\OLId{greek-variable}{Gamma}$ في
كلتا الحالتين.}{}
\end{proof}

\tagprob{FOL}
\begin{prob}
فصّل برهان لمّة الصدق
(\olref[fol][com][mod]{lem:truth})، مع التعديلات التي يحتاج إليها
برهان \olref[fol][com][cpd]{thm:compactness-direct}.
\end{prob}
\tagendprob

\tagprob{notFOL}
\begin{prob}
فصّل برهان لمّة الصدق
(\olref[pl][com][mod]{lem:truth})، مع التعديلات التي يحتاج إليها
برهان \olref[pl][com][cpd]{thm:compactness-direct}.
\end{prob}
\tagendprob

\end{document}
